% !TeX root = ../complete.tex
\section{مدل خطی تأخیر و \lr{Logical Effort} تک‌گیت}
\label{sec:delay-linear-model}

در \hyperref[sec:delay-rc-model]{مدل \lr{RC}}، افزایش ظرفیت بارِ
یک گیت با مقاومت ثابت، تأخیر آن را افزایش می‌دهد. مدل خطی همین
رفتار را با شیب و عرض از مبدأ یک خط بیان می‌کند. این دو پارامتر
امکان مقایسه گیت‌ها را با محاسبه دستی فراهم می‌کنند
\cite[بخش \secnum{۴.۴}، صفحه ۱۵۵]{weste}
\cite[اسلایدهای ۳ تا ۵]{logicaleffortslides}.

\subsection{تأخیر نرمال‌شده و دو سهم آن}
\label{sec:delay-stage-effort}

\(\tau\) مقیاس زمانی مرجع برای یک فناوری و شرایط کاری مشخص است.
در تقریب \lr{RC} کتاب، اگر \(R\) مقاومت مؤثر \lr{Inverter}
واحد در قرارداد زمان نصف تغذیه و \(C\) ظرفیت ورودی یک ترانزیستور
واحد باشد، \(\tau\approx3RC\) است. ظرفیت ورودی خود
\lr{Inverter} مرجع \(3C\) بود. مقدار \(d\) با تقسیم تأخیر
زمانی \(t_{\mathrm{pd}}\) بر \(\tau\) بی‌بعد می‌شود. ظرفیت
ورودی مربوط به گذار را \(C_{\mathrm{in}}\) و ظرفیت خارجی‌ای
را که خروجی راه‌اندازی می‌کند \(C_L\) می‌نامیم.
کتاب و اسلایدها برای این بار خارجی از \(C_{\mathrm{out}}\)
استفاده می‌کنند. در این بخش \(C_L\) را به کار می‌بریم تا با
ظرفیت کل گره خروجی در مدل گذرا اشتباه نشود. روابط مدل چنین‌اند
\cite[روابط \secnum{۴.۲۰} تا \secnum{۴.۲۲}، صفحه ۱۵۵]{weste}
\cite[اسلاید ۴]{logicaleffortslides}.
\begin{equation}
  h=\frac{C_L}{C_{\mathrm{in}}},\qquad
  f=gh,\qquad
  d=\frac{t_{\mathrm{pd}}}{\tau}=f+p=gh+p.
  \label{eq:delay-normalized-gate}
\end{equation}
در این رابطه، \(h\) یا \lr{Electrical Effort} نسبت بار خارجی
به ظرفیت ورودی است. برای چند ورودی مشابه، \(h\) تنها زمانی برابر
تعداد آن‌هاست که ظرفیت هر یک با ظرفیت ورودی گیت راه‌انداز برابر
باشد و ظرفیت سیم قابل صرف‌نظر باشد. \(g\) یا \lr{Logical Effort}
اثر نوع گیت و ورودی موردنظر را بیان می‌کند. \(f\) یا
\lr{Stage Effort} سهم وابسته به بار است. \(p\) یا
\lr{Parasitic Delay} تأخیر نرمال‌شده در \(C_L=0\) است.
ظرفیت‌های درونی که در \(p\) منظور شده‌اند، در \(C_L\) دوباره
شمرده نمی‌شوند.

برای فهم \(g\)، گیت را با یک \lr{Inverter} مقایسه کنید که در
گذار موردنظر همان قدرت راه‌اندازی را دارد. اگر \(C_{\mathrm{in,ref}}\)
ظرفیت ورودی این \lr{Inverter} باشد، تعریف \(g\) چنین است
\cite[بخش \secnum{۴.۴.۱}، صفحه ۱۵۶]{weste}
\cite[اسلاید ۶]{logicaleffortslides}.
\begin{equation}
  g=\frac{C_{\mathrm{in}}}{C_{\mathrm{in,ref}}}.
  \label{eq:delay-gate-logical-effort}
\end{equation}
در تقریب \lr{RC}، مقاومت مؤثر گیت و \lr{Inverter} هم‌قدرت را
برابر و برابر \(R_{\mathrm{eq}}\) می‌گیریم. با تغییر یکنواخت اندازه
\lr{Inverter}، مقاومت و ظرفیت ورودی آن به‌طور معکوس تغییر
می‌کنند و \(R_{\mathrm{eq}}C_{\mathrm{in,ref}}\approx\tau\)
می‌ماند. بنابراین سهم تأخیر بار در واحد \(\tau\) از
\(R_{\mathrm{eq}}C_L/\tau\) به
\(C_L/C_{\mathrm{in,ref}}=gh\) می‌رسد. از این راه روشن می‌شود
که چرا \(g\) شیب نمودار تأخیر نسبت به \(h\) است.

\subsection{عرض ترانزیستورها و ضرایب مدل}
\label{sec:delay-logical-effort-sizing}

فرض کنیم مقاومت یک \lr{pMOS} واحد \(2R\) و مقاومت یک
\lr{nMOS} واحد \(R\) است. در این مدل \lr{Inverter} مرجع
با عرض‌های نسبی ۱ و ۲، در هر دو جهت مقاومت \(R\) و ظرفیت
ورودی \(3C\) دارد. برای \lr{NAND} دوورودی، در مسیر افت دو
\lr{nMOS} سری روشن‌اند. با عرض ۲ برای هرکدام، مقاومت مسیر
\(R/2+R/2=R\) می‌شود. در مسیر صعود، در حالت یک
\lr{pMOS} روشن، عرض ۲ کافی است. هر ورودی به یک
\lr{nMOS} و یک \lr{pMOS} با عرض ۲ وصل است، پس ظرفیت
هر ورودی \(4C\) و \(g=4/3\) است.

در \lr{NOR} دوورودی، دو \lr{pMOS} سری در مسیر صعود
قرار دارند. برای هرکدام عرض ۴ می‌گیریم تا
\(2R/4+2R/4=R\) شود. در مسیر افت، یک \lr{nMOS} با عرض ۱
برای مقاومت \(R\) کافی است. این بار هر ورودی ظرفیت \(5C\)
دارد و \(g=5/3\) می‌شود. با همین استدلال، برای \(n\) ورودی،
عرض ترانزیستورهای سری در \lr{NAND} برابر \(n\) و در
\lr{NOR} برابر \(2n\) است. جدول \ref{tab:delay-simple-gates}
ظرفیت هر ورودی و پارامترهای مدل را برای این گیت‌ها نشان می‌دهد
\cite[شکل \secnum{۴.۲۲} و جدول‌های \secnum{۴.۲}
و \secnum{۴.۳}، صفحات ۱۵۶ و ۱۵۷]{weste}
\cite[اسلایدهای ۶ تا ۹]{logicaleffortslides}.

\begin{table}[htbp]
  \centering
  \caption{پارامترهای گیت‌ها برای مقاومت راه‌اندازی یکسان.
  \(C\) ظرفیت ورودی ترانزیستور واحد و \(p_{\mathrm{inv}}\)
  تأخیر \lr{Parasitic} نرمال‌شده \lr{Inverter} است.}
  \label{tab:delay-simple-gates}
  \begin{tabular}{cccc}
    \toprule
    گیت & ظرفیت هر ورودی & \(g\) & \(p/p_{\mathrm{inv}}\) \\
    \midrule
    \lr{Inverter} & \(3C\) & \(1\) & \(1\) \\
    \lr{NAND} با \(n\) ورودی & \((n+2)C\) & \((n+2)/3\) & \(n\) \\
    \lr{NOR} با \(n\) ورودی & \((2n+1)C\) & \((2n+1)/3\) & \(n\) \\
    \bottomrule
  \end{tabular}
\end{table}

برای ستون آخر جدول، فقط ظرفیت‌های \lr{Diffusion} متصل به
خروجی را می‌شماریم. در \lr{NAND}، یک \lr{nMOS} با عرض
\(n\) و \(n\) ترانزیستور \lr{pMOS} با عرض ۲ به خروجی
متصل‌اند. مجموع عرض‌ها \(3n\) است. در \lr{NOR}، یک
\lr{pMOS} با عرض \(2n\) و \(n\) ترانزیستور \lr{nMOS}
با عرض ۱ باز هم مجموع \(3n\) را می‌دهند. با مقاومت مؤثر برابر
و با فرض ناحیه‌های \lr{Drain} مجزا، تأخیر بی‌بار تقریباً
\(n\) برابر مقدار \lr{Inverter} است. در تقریب \(3C\) برای
ظرفیت \lr{Diffusion} خروجی \lr{Inverter} واحد،
\(p_{\mathrm{inv}}\approx1\) است.

با بزرگ‌کردن یکنواخت یک گیت به اندازه \(k\)، مقاومت به
\(1/k\) مقدار قبلی می‌رسد و ظرفیت داخلی تقریباً \(k\) برابر
می‌شود. از این رو، حاصل‌ضرب آن‌ها و در نتیجه \(p\) در تقریب
اول تقریباً ثابت می‌ماند. گره‌های داخلی می‌توانند این تخمین را
تغییر دهند. برای نمونه، در گذار افت \lr{NAND} کتاب، ترانزیستور
پایینی آخر از همه روشن می‌شود و ظرفیت گره‌های داخلیِ از پیش
شارژشده نیز تخلیه می‌شود. با شماره‌گذاری از سمت زمین،
گره داخلی \(i\) ظرفیت \(nC\) دارد و با خروجی در مقاومت
\(iR/n\) مشترک است. با روش \lr{Elmore}، جمع
سهم خروجی و گره‌های داخلی چنین می‌شود
\cite[شکل \secnum{۴.۲۳} و رابطه \secnum{۴.۲۳}، صفحات ۱۵۷ و ۱۵۸]{weste}.
\[
  t_{\mathrm{pd}}\approx 3nRC+\sum_{i=1}^{n-1}\frac{iR}{n}(nC)
    =\left(3n+\frac{n(n-1)}{2}\right)RC.
\]
جمله دوم نشان می‌دهد چرا تخمین مبتنی بر ظرفیت خروجی به‌تنهایی،
برای پشته‌های سری بلند سهم مهمی از تأخیر را از دست می‌دهد.

\subsection{خواندن نمودار و محاسبه تأخیر \lr{FO4}}
\label{sec:delay-fo4}

در شکل \ref{fig:delay-gate-comparison}، عرض از مبدأ برابر \(p\)
و شیب برابر \(g\) است. با فرض \(p_{\mathrm{inv}}=1\)،
در \(h=3\) به‌ترتیب \(d=4\)، \(d=6\) و \(d=7\) برای
\lr{Inverter}، \lr{NAND} دوورودی و \lr{NOR} دوورودی
به دست می‌آید. مقایسه در \(h\) یکسان، یعنی مقایسه با نسبت بار
به ظرفیت ورودی یکسان. اگر ظرفیت ورودی گیت‌ها متفاوت باشد،
بارهای خارجیِ متناظر نیز متفاوت‌اند.

\begin{figure}[htbp]
  \centering
  \begin{latin}
  \begin{tikzpicture}[x=2.15cm,y=0.47cm]
    \draw[->] (0,0) -- (4.65,0) node[right] {\(h\)};
    \draw[->] (0,0) -- (0,9.05) node[above] {\(d\)};
    \node[below left] at (0,0) {\rl{۰}};
    \foreach \x/\ticklabel in {1/۱,2/۲,3/۳,4/۴} {
      \draw (\x,0.09) -- (\x,-0.09) node[below] {\rl{\ticklabel}};
    }
    \foreach \y/\ticklabel in {1/۱,2/۲,4/۴,6/۶,8/۸} {
      \draw (0.04,\y) -- (-0.04,\y) node[left] {\rl{\ticklabel}};
    }
    \draw[gray,densely dotted] (3,0) -- (3,7);
    \draw[thick] (0,1) -- (3.65,4.65) node[right] {\small Inverter};
    \draw[thick,blue!70!black,dashed] (0,2) -- (3.65,{2+3.65*4/3})
      node[right] {\small NAND2};
    \draw[thick,red!65!black,dashdotted] (0,2) -- (3.65,{2+3.65*5/3})
      node[right] {\small NOR2};
    \fill (3,4) circle (1.6pt);
    \fill[blue!70!black] (3,6) circle (1.6pt);
    \fill[red!65!black] (3,7) circle (1.6pt);
  \end{tikzpicture}
  \end{latin}
  \caption{تأخیر نرمال‌شده \(d\) در برابر \lr{Electrical Effort} بی‌بعد
  \(h\) برای سه گیت، با فرض \(p_{\mathrm{inv}}=1\).
  نمودار از رابطه \eqref{eq:delay-normalized-gate} و ضرایب
  جدول \ref{tab:delay-simple-gates} رسم شده است.}
  \label{fig:delay-gate-comparison}
\end{figure}

در \lr{Fanout-of-4 (FO4)}، یک \lr{Inverter} چهار ورودی
\lr{Inverter} هم‌اندازه را راه‌اندازی می‌کند. در مدل ساده،
\(g=1\)، \(h=4\) و \(p=1\) است. در نمونه کتاب برای یک فناوری
\qty{65}{\nano\metre} با
\(\tau=\qty{3}{\pico\second}\)، تأخیر زمانی چنین به دست می‌آید
\cite[مثال \secnum{۴.۱۰}، صفحه ۱۵۸]{weste}
\cite[اسلاید ۱۲]{logicaleffortslides}.
\[
  d_{\mathrm{FO4}}=1\cdot4+1=5,\qquad
  t_{\mathrm{FO4}}=d_{\mathrm{FO4}}\tau=\qty{15}{\pico\second}.
\]
مقدار ۵ نتیجه مدل نرمال‌شده است و مقدار زمانی به \(\tau\)
وابسته است. به همین دلیل، مقایسه برحسب \lr{FO4} باید با
مرجع همان فناوری و شرایط کاری انجام شود.

\subsection{دوره و فرکانس \lr{Ring Oscillator}}
\label{sec:delay-ring-oscillator}

در حلقه‌ای با تعداد فرد \(N\) از \lr{Inverter}ها، گذار پس
از یک دور با قطبیت مخالف و پس از دور دوم با قطبیت اولیه به
همان گره برمی‌گردد. هر گیت طی یک دوره، یک گذار صعودی و یک
گذار نزولی دارد. اگر گیت‌ها هم‌اندازه باشند و تأخیر هر گذار
به‌ترتیب \(t_\uparrow\) و \(t_\downarrow\) باشد، دوره \(T\) و
فرکانس \(f_{\mathrm{osc}}\) چنین‌اند
\cite[شکل \secnum{۴.۲۴} و مثال \secnum{۴.۱۱}، صفحات ۱۵۸ و ۱۵۹]{weste}
\cite[اسلایدهای ۱۰ و ۱۱]{logicaleffortslides}.
\begin{equation}
  T=N(t_\uparrow+t_\downarrow)=2N\bar t_{\mathrm{pd}},\qquad
  f_{\mathrm{osc}}=\frac{1}{T}.
  \label{eq:delay-ring-frequency}
\end{equation}
در این رابطه، \(\bar t_{\mathrm{pd}}\) میانگین تأخیر صعود و افت
یک مرحله است. اگر بار اضافی و ظرفیت سیم ناچیز باشند، هر
\lr{Inverter} یک ورودی هم‌اندازه را راه‌اندازی می‌کند. بنابراین
\(h=1\) و با فرض \(p=1\)، مقدار \(d=1\cdot1+1=2\) است.
با \(N=31\) و \(\tau=\qty{3}{\pico\second}\)، تأخیر هر
مرحله \(\bar t_{\mathrm{pd}}=\qty{6}{\pico\second}\)، دوره
\(T=\qty{372}{\pico\second}\) و فرکانس
\(f_{\mathrm{osc}}\approx\qty{2.7}{\giga\hertz}\) می‌شود.
عامل ۲ در دوره از دو دور لازم برای بازگشت به قطبیت اولیه می‌آید.
برای محاسبه فرکانس برحسب هرتز، دوره باید برحسب ثانیه باشد.

\subsection{قدرت راه‌اندازی، استخراج ضرایب و حدود مدل}
\label{sec:delay-characterization}

در کتابخانه \lr{Standard Cell}، برچسب‌هایی مثل \lr{X1} و
\lr{X4} اندازه راه‌اندازی را نسبت به یک مرجع کتابخانه بیان
می‌کنند. در مدل آموزشی، اگر \(C_{\mathrm{in,inv}}\) ظرفیت
ورودی \lr{Inverter} واحد باشد، قدرت راه‌اندازی نسبی \(x\)
با نسبت ظرفیت \lr{Inverter} هم‌قدرت به ظرفیت
\lr{Inverter} واحد تعریف می‌شود. با رابطه
\eqref{eq:delay-gate-logical-effort}، این تعریف و صورت دیگر
تأخیر چنین‌اند \cite[روابط \secnum{۴.۲۴} و
\secnum{۴.۲۵}، صفحه ۱۵۹]{weste}.
\begin{equation}
  x=\frac{C_{\mathrm{in}}}{gC_{\mathrm{in,inv}}},\qquad
  d=\frac{C_L/C_{\mathrm{in,inv}}}{x}+p.
  \label{eq:delay-gate-drive}
\end{equation}
برای نمونه، \lr{NAND} دوورودیِ طراحی‌شده با عرض‌های جدول،
ظرفیت هر ورودی \(4C\) و \(g=4/3\) دارد. در نتیجه با مرجع
\(C_{\mathrm{in,inv}}=3C\)، قدرت راه‌اندازی آن \(x=1\) است.
این برابری نشان می‌دهد که ظرفیت ورودی گیت‌های هم‌قدرت می‌تواند
متفاوت باشد.

برای استخراج پارامترها از داده اندازه‌گیری یا \lr{Datasheet}،
در شرایط یکسان خط \(t_{\mathrm{pd}}=a+b h\) را برازش می‌کنیم.
در این رابطه \(a\) عرض از مبدأ و \(b\) شیب، هر دو با واحد زمان‌اند.
اگر داده به شکل تأخیر برحسب ظرفیت بار ارائه شود، شیب زمانی
برحسب \(h\) با ضرب شیب تأخیر برحسب ظرفیت در \(C_{\mathrm{in}}\)
به دست می‌آید. برای \lr{Inverter}، \(g=1\) است و
\(\tau=b_{\mathrm{inv}}\). برای هر گیت، \(g=b/\tau\) و
\(p=a/\tau\) است. در مثال کتاب، شیب \lr{Inverter} واحد
\qty{12.4}{\pico\second} است. برای یکی از ورودی‌های
\lr{NAND} دوورودی، شیب \qty{15.5}{\pico\second} و عرض
از مبدأ \qty{25.4}{\pico\second} است و به ترتیب \(g\approx1.25\)
و \(p\approx2.05\) به دست می‌آید
\cite[بخش \secnum{۴.۴.۵} و روابط \secnum{۴.۲۶}
و \secnum{۴.۲۷}، صفحات ۱۵۹ و ۱۶۰]{weste}.

ضرایب مدل خطی برای بازه مشخصی از بار، شیب ورودی و شرایط کاری
مناسب‌اند. اگر ورودی کندتر شود، ترانزیستور دیرتر قدرت راه‌اندازی
کامل پیدا می‌کند. اگر چند ورودی هم‌زمان تغییر کنند، مسیرهای سری
و موازی می‌توانند تأخیر متفاوتی نسبت به گذار تک‌ورودی داشته باشند.
همچنین \lr{Velocity Saturation} و تغییر ولتاژ تغذیه و آستانه
بر تخمین تأخیر اثر می‌گذارند. در پشته سری، تغییر ولتاژ
\lr{Source} جریان لازم برای شارژ ظرفیت \(C_{gs}\) بین
\lr{Gate} و \lr{Source} را تغییر می‌دهد. ظرفیت \(C_{gd}\)
بین \lr{Gate} و \lr{Drain} نیز ورودی و خروجی را به هم
کوپل می‌کند. در نتیجه، تغییر ولتاژ یکی می‌تواند بار لازم برای
تغییر دیگری را افزایش دهد. هنگام استفاده از ضرایب
جدولی باید گذار فعال و فرض‌های اندازه‌گذاری آن‌ها را مشخص کرد
\cite[بخش \secnum{۴.۴.۶}، صفحات ۱۶۰ تا ۱۶۳]{weste}.
